\documentclass[11pt,a4paper]{article}

\usepackage{kotex}
\usepackage{fontspec}
\setmainfont{Times New Roman}
\setmainhangulfont{AppleMyungjo}
\setsanshangulfont{Apple SD Gothic Neo}

\usepackage[a4paper,top=18mm,bottom=18mm,left=20mm,right=20mm]{geometry}
\usepackage{fancyhdr}
\setlength{\headheight}{24pt}
\usepackage{enumitem}
\usepackage{mdframed}
\usepackage{array}
\usepackage{booktabs}

\setlength{\parindent}{1em}
\linespread{1.18}

\pagestyle{fancy}
\fancyhf{}
\renewcommand{\headrulewidth}{0.6pt}
\fancyhead[L]{\sffamily\bfseries 2026 고1 영어 변형 --- 지문 21}
\fancyhead[R]{\sffamily [shop for words]}
\renewcommand{\footrulewidth}{0pt}
\fancyfoot[C]{\framebox[16mm][c]{\thepage}}

\newcommand{\qno}[1]{\par\vspace{8pt}\noindent{\sffamily\bfseries #1.}\ }
\newcommand{\instr}[1]{{\sffamily #1}\par\vspace{3pt}}
\newlist{choices}{enumerate}{1}
\setlist[choices]{label=,leftmargin=1.5em,itemsep=2pt,topsep=3pt,parsep=0pt}
\newmdenv[linewidth=0.6pt,roundcorner=3pt,innertopmargin=7pt,innerbottommargin=7pt,
          innerleftmargin=9pt,innerrightmargin=9pt,skipabove=5pt,skipbelow=5pt]{pbox}
\newcommand{\nemo}[1]{\fbox{\,#1\,}}

\begin{document}

\begin{center}
{\fontsize{15}{17}\selectfont\sffamily\bfseries [지문 21] 다음 글을 읽고 물음에 답하시오.}
\end{center}
\vspace{4pt}

\begin{pbox}
A classical ideal of truth and knowledge is the concept of a market for ideas: as thinking
individuals, we encounter competing ideas, we evaluate them for their best \nemo{(A) fit / fitting}
to reality, and then let the best and brightest facts \nemo{(B) form / to form} the basis of our
beliefs. That would be the scientist's way, in theory at least. But in practice, most of the time
we are in a market for \rule{3.5cm}{0.4pt}. We shop for words not because they contain ideas,
but because they contain stories about ideas. Often we are not seeking statements as facts to help
us figure out what we should believe. Usually we already know what we believe. What we seek are
statements to justify those beliefs --- the lawyer's way. Language might be good for this function
but we can use it to do \nemo{(C) good / better}.\ By treating language with greater respect, we
can achieve more than simple self-defense.
\end{pbox}

\bigskip

\qno{1}\instr{다음 빈칸에 들어갈 말로 가장 적절한 것은? [3점]}
\begin{choices}
\item ① inspirations
\item ② innovations
\item ③ rationalizations
\item ④ contradictions
\item ⑤ observations
\end{choices}

\bigskip

\qno{2}\instr{(A), (B), (C)의 각 네모 안에서 어법에 맞는 표현으로 가장 적절한 것은?}

\vspace{4pt}
\begin{center}
\begin{tabular}{ccc}
\toprule
\textbf{(A)} & \textbf{(B)} & \textbf{(C)} \\
\midrule
① fitting & form & better \\
② fit & form & better \\
③ fit & to form & better \\
④ fitting & to form & good \\
⑤ fit & form & good \\
\bottomrule
\end{tabular}
\end{center}

\bigskip

\qno{3}\instr{[서술형] 다음 문장을 우리말로 해석하시오.}

\vspace{4pt}
\noindent\textit{By treating language with greater respect, we can achieve more than simple self-defense.}

\vspace{8pt}
\noindent $\Rightarrow$ \rule{\linewidth}{0.4pt}

\vspace{6pt}
\noindent \rule{\linewidth}{0.4pt}

% ===================== Q4 =====================
\qno{4}\instr{윗글의 제목으로 가장 적절한 것은?}
\begin{choices}
\item ① Why Scientific Facts Always Win in Public Debate
\item ② Language as a Weapon for Defending What We Already Believe
\item ③ The Marketplace Where Every Idea Is Judged Fairly
\item ④ How Stories Prevent People from Forming Any Beliefs
\item ⑤ The Decline of Language in Modern Scientific Research
\end{choices}

\end{document}
